\scanpage{30}\mspage{29}
\uline{Cont 4} Si $F,G,$ \deleted{$H$ sont des figures}\inserted{sont deux} figures, telles que tout multistrate de $F$ est comp. à tout multistrate de $G^{*}$, alors $F$ et $G$ sont compatibles (l'inverse étant évident)

% slanted
\margin{\uline{NB} Alors $F$ et $G$ sont ens. comp.}

\deleted{\uline{Cor} Si $F,G,H$ sont \unsure{. tout|. trois} $(\unsure{F_i|F_j}]$}

\deleted{Équi} Ceci implique

\uline{Cor} Soit $(F_i)_{i\in I}$ une famille \deleted{\unsure{finie}} de figures. Pour que \deleted{les $F_i$ soient} cette famille soit une configuration, il f et s que pour tout $i\ne j$ et tout multistrate $\widehat{X}\deleted{\in}$ de $F_i$, $\widehat{Y}$ de $F_j$, $\widehat{X}$ et $\widehat{Y}$ soient compatibles. Si $I$ est fini, cela implique que $F=\bigcup F_i$ est une figure, et \inserted{les} $F_i$ \deleted{est} sont des sous-figures.

\deleted{\uline{NB} Inversement, la validité de cet énoncé pour \deleted{des} tout paire $\{F,G\}$ de figures (savoir, \deleted{si les} le critère de}

\uline{Contrées \deleted{ensemblistes}\inserted{de points}} La contrée est dite ``ensembliste'', ou ``\uline{contrée de points}'' (\unsure{càd|ici|là|i.e.} ses lieux prennent le nom de \uline{points}) si deux figures ensemblistement compatibles sont compatibles (ce qui rend donc tautologique Cont 4) 
